\begin{convention}\label{fga-con-6-0-1}\dcref{6.0.1}\source{fga-explained:page-0153}\uses{}
Fix a field $k$.  Unless stated otherwise, schemes in this chapter are locally
of finite type over $k$ and a point means a $k$-valued point.  Infinitesimal
neighbourhoods of a point are tested by morphisms $\operatorname{Spec}A\to Y$
from local Artinian $k$-algebras sending the closed point to the given point.
In particular, deformations of a closed subscheme $Z\subset X$ are flat families
over $\operatorname{Spec}A$ with special fibre $Z$.
\end{convention}

\begin{convention}\label{fga-con-6-1-1}\dcref{6.1.1}\source{fga-explained:page-0153}\uses{}
For a local ring $A$, write $\mathfrak m_A$ for its maximal ideal.
\end{convention}

\begin{remark}\label{fga-rem-6-1-3}\dcref{6.1.3}\source{fga-explained:page-0153}\uses{}
For a local $k$-algebra $A$ with residue field $k$, the following are equivalent:
\begin{enumerate}
\item $A$ is Artinian;
\item $A$ is finite-dimensional over $k$;
\item $\mathfrak m_A$ is finitely generated and nilpotent.
\end{enumerate}
Thus $A\in(\operatorname{Art}/k)$ exactly when $\operatorname{Spec}A$ is a finite
type $k$-scheme whose reduction is $\operatorname{Spec}k$.
\end{remark}

\begin{remark}\label{fga-rem-6-1-8}\dcref{6.1.8}\source{fga-explained:page-0154}\uses{fga-def-6-1-2}
For a scheme $X$ and a point $p$, the functor of infinitesimal neighbourhoods is
\[
D_{X,p}(A)=\{\operatorname{Spec}A\to X\mid\text{the closed point maps to }p\}.
\]
For a deformation functor $D$, one studies the image of $D(B)\to D(A)$ and
the fibres of this map for every surjection $B\twoheadrightarrow A$ in
$(\operatorname{Art}/k)$.
\end{remark}

\begin{remark}\label{fga-rem-6-1-10}\dcref{6.1.10}\source{fga-explained:page-0154}\uses{}
Every surjection of local Artinian $k$-algebras is a finite composite of small
extensions (extensions with square-zero kernel).  It therefore suffices to
analyze deformation functors on small extensions.
\end{remark}

\begin{remark}\label{fga-rem-6-1-11}\dcref{6.1.11}\source{fga-explained:page-0154}\uses{fga-def-6-1-2}
For a deformation functor $D$, the map $D(B)\to D(k)$ is surjective.  Its
distinguished element is the image of the unique element of $D(k)$ under the
structure map $k\to B$.
\end{remark}

\begin{definition}\label{fga-def-6-1-12}\dcref{6.1.12}\source{fga-explained:page-0155}\uses{}
Let $(\operatorname{Loc}/k)$ be the category of local $k$-algebras $R$ with
residue field $k$ and finite-dimensional tangent space
$T_R=(\mathfrak m_R/\mathfrak m_R^2)^\vee$.  Set
\[
h_R(A)=\operatorname{Hom}_{k\text{-alg}}(R,A),\qquad
\widehat R=\varprojlim_nR/\mathfrak m_R^n.
\]
Then $h_R$ is a deformation functor; $R$ is complete if $R\simeq\widehat R$.
Write $(\operatorname{CLoc}/k)$ for the full subcategory of complete rings.
\end{definition}

\begin{remark}\label{fga-rem-6-1-14}\dcref{6.1.14}\source{fga-explained:page-0155}\uses{fga-def-6-1-12}
For $R=\widehat{\mathcal O}_{X,p}$, $h_R(A)$ is naturally identified with
$D_{X,p}(A)$.
\end{remark}

\begin{remark}\label{fga-rem-6-1-15}\dcref{6.1.15}\source{fga-explained:page-0155}\uses{fga-def-6-1-16}
If $R=\widehat{\mathcal O}_{X,p}$, then $\dim_kT_R$ is the least dimension of a
smooth scheme having an open neighbourhood of $p$ embedded as a closed
subscheme.
\end{remark}

\begin{definition}\label{fga-def-6-1-16}\dcref{6.1.16}\source{fga-explained:page-0155}\uses{}
The dimension of $T_R$ is the embedding dimension of $R$.
\end{definition}

\begin{remark}\label{fga-rem-6-1-20}\dcref{6.1.20}\source{fga-explained:page-0156}\uses{fga-thm-6-1-19,fga-def-6-1-21}
For a small extension $0\to M\to B\to A\to0$, a tangent-obstruction sequence
\[
0\to T_D\otimes_kM\to D(B)\to D(A)
 \xrightarrow{\operatorname{ob}}T^2_D\otimes_kM
\]
means that $\operatorname{ob}(a)=0$ exactly when $a$ lifts, and the set of
liftings is a torsor under $T_D\otimes_kM$.  A morphism of small extensions
induces a commutative diagram of these sequences, with equivariant maps on the
torsors.
\end{remark}
\begin{proof}
For $R=S/J$ with $S=k[[t_1,\ldots,t_d]]$ and $J\subset\mathfrak m_S^2$, choose
a lift $S\to B$ of a map $R\to A$.  Two lifts differ by a derivation
$S\to M$, hence by an element of $(\mathfrak m_S/\mathfrak m_S^2)^\vee\otimes_kM$.
The difference vanishes on $J$ precisely when the obstruction map
$J/\mathfrak m_SJ\to M$ is zero.  This proves exactness and functoriality.
\end{proof}

\begin{remark}\label{fga-rem-6-1-22}\dcref{6.1.22}\source{fga-explained:page-0157}\uses{fga-def-6-1-21}
If $T^2_D\hookrightarrow T'^2_D$ is injective, replacing the obstruction map by
its composite gives another tangent-obstruction theory.  The tangent space is
intrinsic, whereas the obstruction space is not uniquely determined.
\end{remark}

\begin{proposition}\label{fga-pro-6-1-23}\dcref{6.1.23}\source{fga-explained:page-0157}\uses{fga-def-6-1-9,fga-def-6-1-21}
For a functor with a generalized tangent-obstruction theory, its tangent vector
space is determined canonically.
\end{proposition}
\begin{proof}
Put $A_1=k[\varepsilon]/(\varepsilon^2)$.  The distinguished element of
$D(A_1)$ identifies $D(A_1)$ with the tangent torsor, hence with $T_D$.
The maps $A_2=k[\varepsilon_1,\varepsilon_2]/(\varepsilon_1,\varepsilon_2)^2$
$\to A_1$ and $\varepsilon_1+\varepsilon_2\mapsto\varepsilon$ transport the
product of the two fibres to an addition on $D(A_1)$.  Functoriality makes scalar
multiplication and addition agree with the given vector-space structure, so the
identification is canonical.
\end{proof}

\begin{remark}\label{fga-rem-6-2-2}\dcref{6.2.2}\source{fga-explained:page-0158}\uses{fga-def-6-2-1}
Here ``pro-representable'' means represented by $h_R$ for a complete local
$k$-algebra $R$ with finite-dimensional tangent space.  This is slightly
stronger than the categorical notion of a pro-object in $(\operatorname{Art}/k)$.
\end{remark}

\begin{remark}\label{fga-rem-6-2-3}\dcref{6.2.3}\source{fga-explained:page-0158}\uses{fga-thm-6-1-19}
If $D=h_R$, choose $R\simeq S/J$ with $S=k[[t_1,\ldots,t_d]]$ and
$J\subset\mathfrak m_S^2$.  Then the tangent space is
$(\mathfrak m_S/\mathfrak m_S^2)^\vee$ and the canonical obstruction space is
$(J/\mathfrak m_SJ)^\vee$.
\end{remark}

\begin{remark}\label{fga-rem-6-3-2}\dcref{6.3.2}\source{fga-explained:page-0160}\uses{fga-pro-6-1-23,fga-def-6-3-1}
A map $h_R\to F$ is a hull exactly when it is smooth and induces an isomorphism
on tangent spaces.  The formal criterion for smoothness identifies this with
formal smoothness of the corresponding morphism of schemes; a formally smooth
morphism inducing an isomorphism on tangent spaces is formally etale.
\end{remark}

\begin{lemma}\label{fga-lem-6-3-3}\dcref{6.3.3}\source{fga-explained:page-0160}\uses{fga-def-6-3-1,fga-pro-6-1-23}
If $h_R\to F$ is a pro-representable hull and $F$ has a tangent-obstruction
theory $(T^1,T^2)$, then $(T^1,T^2)$ is also a tangent-obstruction theory for
$h_R$.  In particular $T^1\simeq T_R$ and the obstruction dimensions of $R$
are bounded by those of $F$.
\end{lemma}
\begin{proof}
Smoothness of the hull identifies liftings of $h_R$ with liftings in $F$ and
preserves their torsor actions.  The tangent identification is Proposition
\ref{fga-pro-6-1-23}; the obstruction map is transported through the hull.
\end{proof}

\begin{remark}\label{fga-rem-6-3-6}\dcref{6.3.6}\source{fga-explained:page-0161}\uses{fga-def-6-3-1}
Two pro-representable hulls of a functor are isomorphic over that functor.  The
isomorphism need not be unique because automorphisms act trivially on the tangent
space but may act nontrivially on higher Artin quotients.
\end{remark}

\begin{lemma}\label{fga-lem-6-4-3}\dcref{6.4.3}\source{fga-explained:page-0163}\uses{}
Let $B\twoheadrightarrow A$ be a small extension of rings with square-zero
kernel $M$, and let $\mathcal Q'$ be coherent on $X_B$.  Then $\mathcal Q'$ is
flat over $B$ if and only if
\[
\mathcal Q'\otimes_BM\longrightarrow\mathcal Q'
\quad\text{is injective, and}\quad
\mathcal Q'\otimes_BA\text{ is flat over }A.
\]
\end{lemma}
\begin{proof}
Flatness over $B$ implies both assertions by the local criterion for flatness.
Conversely, reduce to a small extension and test injectivity of
$\mathcal Q'\otimes_BI\to\mathcal Q'$ for every ideal $I\subset B$.  Apply the
exact sequence $0\to M\cap I\to I\to I/(M\cap I)\to0$ and tensor with
$\mathcal Q'$.  Flatness over $A$ gives exactness for the quotient, while the
assumed injectivity handles $M\cap I$; hence the map is injective for every
$I$, and the local criterion proves $B$-flatness.
\end{proof}

\begin{lemma}\label{fga-lem-6-4-4}\dcref{6.4.4}\source{fga-explained:page-0163}\uses{}
For an exact sequence of coherent sheaves
\[
0\longrightarrow\mathcal E\longrightarrow\mathcal F
 \xrightarrow{\pi}\mathcal Q\longrightarrow0,
\]
$\pi$ has a section if and only if its class in
$\operatorname{Ext}^1(\mathcal Q,\mathcal E)$ vanishes.  When nonempty, the set
of sections is a torsor under $\operatorname{Hom}(\mathcal Q,\mathcal E)$.
\end{lemma}
\begin{proof}
The sequence represents its Yoneda extension class.  A splitting is exactly a
zero class.  The difference of two splittings factors uniquely through
$\mathcal Q\to\mathcal E$, and adding such a homomorphism to a splitting gives
all splittings, proving the torsor assertion.
\end{proof}

\begin{remark}\label{fga-rem-6-4-6}\dcref{6.4.6}\source{fga-explained:page-0165}\uses{}
In the extension problem, the space
$\operatorname{Hom}_{X_A}(\mathcal E,\mathcal Q\otimes_A M)$ may equivalently
be written $\operatorname{Hom}_{X_k}(\mathcal E\otimes_Ak,
\mathcal Q\otimes_A M)$.
\end{remark}

\begin{corollary}\label{fga-cor-6-4-10}\dcref{6.4.10}\source{fga-explained:page-0166}\uses{fga-thm-6-4-5}
For the Hilbert functor of a closed subscheme $Z\subset X$, a generalized
tangent-obstruction theory is
\[
T^1=\operatorname{Hom}_X(\mathcal I_Z,\mathcal O_Z),\qquad
T^2=\operatorname{Ext}^1_X(\mathcal I_Z,\mathcal O_Z).
\]
\end{corollary}
\begin{proof}
Apply the extension criterion to $0\to\mathcal I_Z\to\mathcal O_X\to
\mathcal O_Z\to0$; first-order changes are homomorphisms
$\mathcal I_Z\to\mathcal O_Z$, and the Cech obstruction to gluing lies in the
displayed Ext group.
\end{proof}

\begin{corollary}\label{fga-cor-6-4-11}\dcref{6.4.11}\source{fga-explained:page-0166}\uses{fga-cor-6-4-10,fga-cor-6-2-6}
Let $X$ be quasiprojective and $[Z]\in\operatorname{Hilb}(X)$ with $Z$ proper.
Set
\[
d=\dim_k\operatorname{Hom}_X(\mathcal I_Z,\mathcal O_Z),\qquad
r=\dim_k\operatorname{Ext}^1_X(\mathcal I_Z,\mathcal O_Z).
\]
Then
\[
d\ge\dim_{[Z]}\operatorname{Hilb}(X)\ge d-r.
\]
Equality with $d$ implies that the Hilbert scheme is nonsingular at $[Z]$;
equality with $d-r$ implies it is a local complete intersection there.
\end{corollary}
\begin{proof}
The tangent-obstruction sequence gives the upper bound by the tangent dimension
and the lower bound by the number of independent equations.  Vanishing of the
obstruction map gives smoothness, while equality with the expected codimension
$r$ gives a regular sequence in the completed local ring, hence a local complete
intersection.
\end{proof}

\begin{theorem}\label{fga-the-6-4-12}\dcref{6.4.12}\source{fga-explained:page-0166}\uses{fga-def-6-1-21}
Let $X$ be separated and $f:X\to Y$ a morphism with $Y$ nonsingular.  The
deformation functor $D_f$ of $f$ has tangent space
\[
T^1=H^0(X,f^*T_Y)
\]
and obstruction space $T^2=H^1(X,f^*T_Y)$; for every small extension its
tangent-obstruction sequence is exact.
\end{theorem}
\begin{proof}
On affine opens, a lifting across a square-zero extension differs from another
by a derivation into the kernel, hence by a section of $f^*T_Y$; affine liftings
are unobstructed.  Choose affine covers of $X$ subordinate to affine opens of
$Y$.  Differences of local liftings on pairwise overlaps form a Cech $1$-cocycle
with values in $f^*T_Y$.  Its class is independent of choices, vanishes exactly
when the local liftings glue, and the set of gluings is a torsor under global
sections.  Separation of $X$ makes the overlap comparisons effective.
\end{proof}

\begin{proposition}\label{fga-pro-6-5-2}\dcref{6.5.2}\source{fga-explained:page-0167}\uses{fga-def-6-4-2}
If $Z$ is a local complete intersection in $X$, then its deformation functor has
\[
T^1=H^0(Z,\mathcal N_{Z/X}),\qquad
T^2=H^1(Z,\mathcal N_{Z/X}).
\]
\end{proposition}
\begin{proof}
For an lci immersion, $\mathcal I_Z/\mathcal I_Z^2$ is locally free and
$\mathcal N_{Z/X}=\mathcal H om(\mathcal I_Z/\mathcal I_Z^2,\mathcal O_Z)$.
The affine lci lifting criterion is unobstructed; on an affine cover, differences
of liftings are sections of $\mathcal N_{Z/X}$ and their Cech class is the only
obstruction.  This gives the displayed groups.
\end{proof}

\begin{remark}\label{fga-rem-6-5-3}\dcref{6.5.3}\source{fga-explained:page-0168}\uses{}
If $X$ is separated and smooth, the deformation functor of $X$ has tangent
space $H^1(X,T_X)$ and no obstruction beyond the corresponding Cech
$H^2$-class; affine smooth schemes have no nontrivial infinitesimal deformations.
\end{remark}
\begin{proof}
Trivialize a deformation on an affine open cover.  The transition maps are
infinitesimal automorphisms, represented by sections of $T_X$; their failure to
satisfy the cocycle condition is a Cech $2$-cocycle.  Changing local
trivializations changes the cocycle by a coboundary, yielding the asserted
tangent and obstruction groups.
\end{proof}
